\documentclass[10pt,a4paper]{article}

\usepackage[a4paper,margin=1.05cm]{geometry}
\usepackage[hidelinks,unicode]{hyperref}
\usepackage{enumitem}
\usepackage{tabularx}
\usepackage{xcolor}
\usepackage{array}
\usepackage{graphicx}
\usepackage{fontspec}

\setmainfont{DejaVu Sans}
\definecolor{accent}{HTML}{1F5F8B}
\definecolor{textgray}{HTML}{4B5563}
\hypersetup{
  pdftitle={Aleksander Borodin -- ML Engineer | LLMs \& AI Agents},
  pdfauthor={Aleksander Borodin},
  pdfsubject={Curriculum vitae}
}

\pagestyle{empty}
\setlength{\parindent}{0pt}
\setlength{\parskip}{0pt}
\setlength{\tabcolsep}{3pt}
\emergencystretch=2em
\setlist[itemize]{leftmargin=1.0em,itemsep=0.06em,topsep=0.06em,parsep=0pt,partopsep=0pt}
\renewcommand\labelitemi{--}

\newcommand{\sectiontitle}[1]{%
  \vspace{0.42em}%
  {\normalsize\bfseries\color{accent} #1}\par
  \vspace{0.10em}\hrule\vspace{0.24em}%
}

\newcommand{\entry}[4]{%
  \textbf{#1}\hfill {\color{textgray}#2}\par
  {\color{textgray}#3}\par
  #4\vspace{0.25em}%
}

\newcommand{\achievement}[2]{%
  #1\hfill{\color{textgray}#2}\par
}

\begin{document}
\fontsize{9.25}{10.4}\selectfont
\raggedright

\begin{minipage}[c]{0.78\textwidth}
  {\LARGE\bfseries Aleksander Borodin}\par
  \vspace{0.04em}
  {\large ML Engineer | LLMs \& AI Agents}\par
  \vspace{0.16em}
  \href{mailto:apdorodin765@gmail.com}{apdorodin765@gmail.com} \quad
  +7 908 971 16 97 \quad
  \href{https://github.com/aleksanderborodin}{github.com/aleksanderborodin}\par
  \href{https://t.me/aleksanderbor}{Telegram: @aleksanderbor} \quad
  \href{https://aleksanderborodin.github.io/}{aleksanderborodin.github.io}
\end{minipage}
\hfill
\begin{minipage}[c]{0.15\textwidth}
  \raggedleft
  \setlength{\fboxsep}{0pt}
  \setlength{\fboxrule}{0.6pt}
  \fcolorbox{accent}{white}{\includegraphics[width=2.25cm]{assets/profile-original.jpg}}
\end{minipage}

\sectiontitle{Summary}
ML engineer and Applied Mathematics and Physics student at MIPT. I build LLM systems, from RAG and model integrations to multi-agent workflows and agent-driven evolution: iterative search and improvement of solutions. My experience includes benchmark development, model evaluation, and applied optimization.

\sectiontitle{Education}
\entry{Bachelor's degree, Applied Mathematics and Physics}{2023--2027}{FPMI MIPT, in progress; GPA 7.8/10}{Coursework: mathematical analysis, probability, linear algebra, algorithms, computational complexity, ML, physics, economics.}

\sectiontitle{Professional Experience}
\entry{ML Engineer -- Research Internship}{Feb--Aug 2026 (6 months)}{CKM, MIPT}{%
\begin{itemize}
  \item Developed LLM-based agents and benchmarks to evaluate their performance.
  \item Worked on agent-driven evolution: systems in which agents iteratively generate, evaluate, and improve solutions.
  \item Contributed to two research papers: one under review at ACL and another in preparation for submission.
\end{itemize}
}

\entry{Product / AI Engineer}{Jan--May 2026}{\href{https://modelgate.ru}{ModelGate.ru} -- LLM API gateway and proxy service}{%
\begin{itemize}
  \item Developed an LLM API gateway providing a unified interface to models from multiple providers.
  \item Integrated provider APIs and implemented request routing with model and provider selection.
  \item Built core service components, including the backend and public website. Stack: Python, LLM APIs.
\end{itemize}
}

\entry{Model and Algorithm Developer}{Sep 2024--Feb 2025}{Regional Bank, client project}{%
\begin{itemize}
  \item Developed a cash collection optimization algorithm; historical-data analysis projected a 20M RUB annual profit increase.
  \item Built a CBR web-scraping/analysis pipeline for data unavailable via API and an internal document Q\&A RAG system.
  \item Stack: Python, NumPy, Pandas, RAG, web scraping.
\end{itemize}
}

\sectiontitle{Projects}
\entry{Idea Evolve}{2026}{\href{https://github.com/aleksanderborodin/idea_evolve}{github.com/aleksanderborodin/idea\_evolve}}{%
\begin{itemize}
  \item Evolutionary multi-agent framework: specialized Claude agents work in parallel across generations and build a shared knowledge base.
  \item Improved Sidon set search from 66 to 105 elements in 7 autonomous generations, passing the $\sqrt{N}\approx100$ reference for $N=10000$.
  \item Implemented role-specialized agents, resumable filesystem state, evaluation/knowledge extraction, and a Flask dashboard.
  \item Stack: Python 3.12+, Claude Code CLI, OpenCode, Anthropic API, Flask, multi-agent orchestration.
\end{itemize}
}

\entry{Prefix Injection Jailbreak}{2025}{\href{https://github.com/aleksanderborodin/prefix_injection_jailbreak}{github.com/aleksanderborodin/prefix\_injection\_jailbreak}}{%
\begin{itemize}
  \item Researched Choice-Blindness-inspired prefix attacks; templates made Gemini 2.5 Pro and DeepSeek R1 complete MaliciousInstruct prompts with near-100\% success vs. a near-zero baseline.
  \item Stack: Python, LLM APIs, model evaluation, prompt/context engineering.
\end{itemize}
}

\entry{Yoola Explain}{2025}{\href{https://github.com/aleksanderborodin/yoola_explain}{github.com/aleksanderborodin/yoola\_explain}}{%
\begin{itemize}
  \item Browser extension that extracts Terms of Service and returns structured LLM summaries: key points, data collection notes, and warnings.
  \item Added multi-language support and caching. Stack: JavaScript, FastAPI, OpenRouter, LLM APIs.
\end{itemize}
}

\sectiontitle{Skills}
\begin{tabularx}{\textwidth}{@{}>{\bfseries}p{2.65cm}X@{}}
Programming & Python advanced; C++ intermediate; C and Assembler basic code reading. \\
Data / ML & PyTorch basic; NumPy; Pandas; PostgreSQL; Matplotlib; ClearML; Excel. \\
NLP / LLM & BM25; embeddings; RAG; LLM APIs; evaluation; prompt/context engineering; jailbreak testing. \\
Agentic tools & Claude Code CLI; OpenCode; multi-agent orchestration; filesystem-based workflows. \\
Backend & FastAPI; Flask; LangChain; LangGraph; Linux home server administration and workstation usage. \\
Languages & Russian native; English B2+ certified by MIPT placement test. \\
\end{tabularx}

\sectiontitle{Programs \& Achievements}
\entry{\href{https://airi.net/ru/summer-school-2026/}{Summer with AIRI 2026 -- Artificial Intelligence Summer School}}{Jul 21--Aug 4, 2026}{Final team project: \href{https://github.com/n4rly-boop/ideas-lake}{Ideas Lake} -- long-term memory for evolutionary AI workflows using hybrid retrieval and a knowledge graph.}{}
\achievement{\textbf{Gold Medal}, \href{https://iepho.ru/}{International Experimental Physics Olympiad (IEPhO)}}{2023}
\achievement{\textbf{Prize-winner}, \href{https://vos.olimpiada.ru/}{All-Russian School Olympiad in Physics}, final stage — Grades 10 and 11}{2022, 2023}
\achievement{Winner, team Wiki-races, Sirius}{2023}

\end{document}
